
% -------------------------------------------------------
%  English Abstract
% -------------------------------------------------------

\begin{latin}

\begin{center}
\textbf{Abstract}
\end{center}
\baselineskip=.8\baselineskip

This project designs, implements and evaluates a prototype Persian
speech-to-speech system that operates in full-duplex mode: the user microphone
stays open while the assistant speaks, and the system reacts to interruptions.
Two tracks were pursued: low-rank adaptation of the end-to-end Moshika model to
Persian, and a streaming cascade of NeMo speech recognition, the
Qwen3-4B-Instruct-2507 responder and Piper speech synthesis. A conversational
corpus of 6,754 user/response pairs (108.584 hours, session-disjoint splits) was
derived from an archive of Persian spoken content.

The final cascade detects the end of speech automatically and speculatively,
streams the reply clause by clause, and handles interruptions in two stages: it
ducks the assistant after one frame of speech louder than the expected
residual echo, then confirms lexically with the speech recognizer after
discarding listener tokens and the assistant's own echoed words. In automated
real-browser loopback, end-of-speech to first assistant audio had a mean of
372 ms and a 90th percentile of 431 ms, with every turn under 500 ms and the
GPU side of the stack using 11.1 GiB. Interruption handling reached accuracy 0.933 and F1 0.880, with a
median of 120 ms to the duck, in a deployment simulation built from real
recordings, and accuracy 0.825 (F1 0.817) in the browser; an energy/F0/MFCC classifier
alone reached only 0.574 on real podcast events. The responder passed all seven
predeclared automatic gates of a 40-row final test (relevance 3.150 versus
1.425). All components were re-selected on fresh panels, including speech
recognition on human transcripts and synthesis with two automatic listeners.

The direct track produced no deployable model. Moshi's English tokenizer needs
6.6 tokens per Persian word and exceeds its one-token-per-frame budget;
phoneme romanization doubled the best runtime result but the speech was still
not intelligible, and a speech-input Gemma 4 model also underperformed the
cascade. Latency and interruption evidence is automated and on an H100 (the
memory of a 12 GB card was respected, its compute speed was not), labels are
automatic, and the results of the human listening study are not yet available.

\bigskip\noindent\textbf{Keywords}:
Speech-to-speech, full-duplex dialogue, barge-in detection, Persian language,
large language models, streaming cascade, tokenization, parameter-efficient adaptation

\end{latin}
